% ---------------------------------------------------------------------------
% Preâmbulo comum do livro
% ---------------------------------------------------------------------------

% --- Codificação e idioma ---------------------------------------------------
\usepackage[T1]{fontenc}
\usepackage[utf8]{inputenc}
\usepackage[brazil]{babel}
\usepackage{lmodern}
\usepackage{microtype}

% --- Matemática -------------------------------------------------------------
\usepackage{amsmath}
\usepackage{amssymb}
\usepackage{amsfonts}
\usepackage{mathtools}
\usepackage{bm}

% --- Tabelas e figuras ------------------------------------------------------
\usepackage{graphicx}
\usepackage{epstopdf}
\epstopdfsetup{update,prepend}
\usepackage{booktabs}
\usepackage{longtable}
\usepackage{multirow}
\usepackage{array}
\usepackage{tabularx}
\usepackage{rotating}
\usepackage[font=small,labelfont=bf]{caption}
\usepackage{subcaption}
\usepackage{float}

% --- Cores e caixas ---------------------------------------------------------
\usepackage{xcolor}
\usepackage[most]{tcolorbox}
\tcbuselibrary{skins,breakable}

% --- Listagens de código ----------------------------------------------------
\usepackage{listings}
\definecolor{codfundo}{RGB}{247,247,247}
\definecolor{codcomentario}{RGB}{100,120,100}
\definecolor{codpalavra}{RGB}{20,60,140}
\definecolor{codtexto}{RGB}{150,60,40}
\lstdefinestyle{python}{
  language         = Python,
  basicstyle       = \ttfamily\footnotesize,
  keywordstyle     = \color{codpalavra}\bfseries,
  commentstyle     = \color{codcomentario}\itshape,
  stringstyle      = \color{codtexto},
  numbers          = none,
  showstringspaces = false,
  breaklines       = true,
  breakatwhitespace= true,
  columns          = fullflexible,
  backgroundcolor  = \color{codfundo},
  frame            = single,
  rulecolor        = \color{black!25},
  framesep         = 5pt,
  xleftmargin      = 5pt,
  xrightmargin     = 5pt,
  aboveskip        = 1.1em,
  belowskip        = 1.1em,
  literate         = {á}{{\'a}}1 {ã}{{\~a}}1 {ç}{{\c c}}1 {é}{{\'e}}1
                     {ê}{{\^e}}1 {í}{{\'i}}1 {ó}{{\'o}}1 {õ}{{\~o}}1 {ú}{{\'u}}1
}
\lstset{style=python}

% --- Unidades e números -----------------------------------------------------
\usepackage{siunitx}
\sisetup{
  output-decimal-marker = {,},
  group-separator       = {.},
  group-minimum-digits  = 4,
  detect-weight         = true,
  detect-family         = true
}

% --- Layout -----------------------------------------------------------------
\usepackage[a4paper,inner=3cm,outer=2.5cm,top=3cm,bottom=2.5cm]{geometry}
\usepackage{emptypage}
\usepackage{enumitem}
% Evita o aviso "Negative labelwidth" nas listas de descricao com leftmargin=0pt
\setlist[description]{labelwidth=0pt,labelsep=0pt}
\usepackage{makeidx}
\makeindex

% --- Bibliografia -----------------------------------------------------------
\usepackage[round,authoryear]{natbib}
\bibliographystyle{apalike}
\renewcommand{\bibname}{Referências}

% --- Hyperref e referências cruzadas (carregar por último) ------------------
\usepackage[hidelinks,breaklinks=true]{hyperref}
\hypersetup{
  pdftitle  = {Aprendizado de Mundo Aberto em Poços de Petróleo},
  pdfauthor = {Lucas Gouveia Omena Lopes},
  pdfsubject= {Detecção, classificação e explicação de anomalias em poços de petróleo},
  pdfkeywords = {mundo aberto, detecção de anomalias, poços de petróleo, aprendizado profundo, LLM}
}
\usepackage[brazilian]{cleveref}

% ---------------------------------------------------------------------------
% Nomes em português para o cleveref
% ---------------------------------------------------------------------------
\crefname{chapter}{Capítulo}{Capítulos}
\Crefname{chapter}{Capítulo}{Capítulos}
\crefname{part}{Parte}{Partes}
\Crefname{part}{Parte}{Partes}
\crefname{section}{Seção}{Seções}
\Crefname{section}{Seção}{Seções}
\crefname{figure}{Figura}{Figuras}
\Crefname{figure}{Figura}{Figuras}
\crefname{table}{Tabela}{Tabelas}
\Crefname{table}{Tabela}{Tabelas}
\crefname{equation}{Equação}{Equações}
\Crefname{equation}{Equação}{Equações}

% ---------------------------------------------------------------------------
% Ambientes didáticos
% ---------------------------------------------------------------------------
\definecolor{corDefinicao}{HTML}{1F4E79}
\definecolor{corExemplo}{HTML}{2E7D32}
\definecolor{corDestaque}{HTML}{B35C00}
\definecolor{corExercicio}{HTML}{5B2C6F}
\definecolor{corResultado}{HTML}{8B1A1A}

\newtcolorbox[auto counter,number within=chapter]{definicao}[2][]{%
  breakable, enhanced, colback=corDefinicao!4, colframe=corDefinicao,
  fonttitle=\bfseries, title={Definição~\thetcbcounter\ --- #2}, #1}

\newtcolorbox[auto counter,number within=chapter]{exemplo}[2][]{%
  breakable, enhanced, colback=corExemplo!4, colframe=corExemplo,
  fonttitle=\bfseries, title={Exemplo~\thetcbcounter\ --- #2}, #1}

\newtcolorbox{destaque}[1][Em resumo]{%
  breakable, enhanced, colback=corDestaque!5, colframe=corDestaque,
  fonttitle=\bfseries, title={#1}}

\newtcolorbox{resultadochave}[1][Resultado-chave]{%
  breakable, enhanced, colback=corResultado!4, colframe=corResultado,
  fonttitle=\bfseries, title={#1}}

\newtcolorbox{objetivos}{%
  breakable, enhanced, colback=black!3, colframe=black!55,
  fonttitle=\bfseries, title={Objetivos de aprendizagem}}

\newenvironment{exercicios}%
  {\section*{Exercícios}\addcontentsline{toc}{section}{Exercícios}%
   \begin{enumerate}[label=\textbf{\arabic*.},leftmargin=*,itemsep=0.4em]}%
  {\end{enumerate}}

\newenvironment{leituras}%
  {\section*{Leituras recomendadas}\addcontentsline{toc}{section}{Leituras recomendadas}%
   \begin{itemize}[leftmargin=*,itemsep=0.3em]}%
  {\end{itemize}}

\newcommand{\resumocapitulo}[1]{%
  \section*{Resumo do capítulo}\addcontentsline{toc}{section}{Resumo do capítulo}#1}

% ---------------------------------------------------------------------------
% Macros de notação
% ---------------------------------------------------------------------------
\newcommand{\vect}[1]{\bm{#1}}
\newcommand{\mat}[1]{\bm{#1}}
\newcommand{\R}{\mathbb{R}}
\newcommand{\E}{\mathbb{E}}
\newcommand{\prob}{\mathbb{P}}
\DeclareMathOperator*{\argmax}{arg\,max}
\DeclareMathOperator*{\argmin}{arg\,min}
\DeclareMathOperator{\std}{std}

% Primeira ocorrência de termo estrangeiro
\newcommand{\ing}[1]{\emph{#1}}
% Termo em português seguido do original em inglês
\newcommand{\termo}[2]{#1\index{#1}\ (\ing{#2})}

% Siglas frequentes
\newcommand{\treW}{3W}
